\section{Введение}
\subsection*{Цель} \addcontentsline{toc}{subsection}{Цель}
Целью лабораторной работы является исследование дифракции Фраунгофера на препятствиях в виде одной щели, двух щелей и решётки из 15
щелей.

\subsection*{Задачи} \addcontentsline{toc}{subsection}{Задачи}
\begin{enumerate}
	
	\item Построить теоретические графики зависимости $I(\theta)$ для дифракции на одной, двух и 15 щелях.
	\item Измерить углы на которых наблюдаются минимумы интенсивности при дифракции на одной и двух щелях, и соответствующие углы для максимумов на решётке из 15 щелей.
	\item Наблюдая дифракцию на двух щелях, определить размер источника свет(ширину входной щели коллиматора), при которой наступает размытие дифракционной картины и сравнить с теорией.
	
\end{enumerate}


\subsection*{Приборы и оборудование} \addcontentsline{toc}{subsection}{Приборы и оборудование}
\begin{enumerate}
	\item Гониометр ГС-30
	\item Решётки с одной, двумя и 15 щелями
\end{enumerate}